\section{Harmonic tails, centered moments, and the source matrix}
\label{sec:framework}

For an integer $s\ge2$, write
\begin{equation}
 A_s(n)=\zeta(s)-H_n^{(s)}=\zeta(s,n+1),
 \qquad
 H_n^{(s)}=\sum_{m=1}^n\frac1{m^s},
 \qquad x_n=n+\frac12.
 \label{eq:tail}
\end{equation}
Thus $A_s(n)$ is simultaneously a generalized harmonic tail, a
Hurwitz-zeta value, and a normalized polygamma value~\cite{DLMF}:
\begin{equation}
 A_s(n)=\frac{(-1)^s}{(s-1)!}\psi^{(s-1)}(n+1).
 \label{eq:polygamma-tail}
\end{equation}
For even $s$, the sign in \eqref{eq:polygamma-tail} is positive.

For even $p,q\ge2$, the source family is
\begin{equation}
 D_{p,q}(w)=\sum_{n=0}^{\infty}A_p(n)A_q(n)x_n^{-w},
 \qquad \Re w>3-p-q.
 \label{eq:dirichlet-family}
\end{equation}
The previous reports continue this function meromorphically and prove
that it is regular at the even nonpositive arguments needed below.  We
write
\[
 D_r=D_{p,q}(-2r),\qquad r\ge0.
\]
The centered Bernoulli transform is
\begin{equation}
 C_j=\sum_{r=0}^j
       \binom{2j+1}{2r}2^{-2(j-r)}D_r,
 \qquad j\ge0.
 \label{eq:C-transform}
\end{equation}
Its inverse is
\begin{equation}
 D_r=\sum_{j=0}^r\frac{\binom{2r}{2j}}{2j+1}
 B_{2r-2j}\!\left(\frac12\right)C_j.
 \label{eq:C-inverse}
\end{equation}
Both transforms are triangular with nonzero diagonal.

\subsection{Late coordinates}

Put $p=2a$, $q=2b$, with $1\le a<b$, and introduce formal labels
\[
 Z_0=\zeta(0)=-\frac12,
 \qquad Z_k=\zeta(2k)\quad(1\le k<b).
\]
The source calculation proves that the late centered moments have the
rational coordinate expansion
\begin{equation}
 C_{b+t}=\sum_{k=0}^{b-1}(M_{a,b})_{tk}Z_k,
 \qquad 0\le t<b.
 \label{eq:late-coordinates}
\end{equation}
Set
\[
 \mu_j=
 \begin{cases}
 B_{2j}/(2j)!,&j\ge0,\\
 0,&j<0.
 \end{cases}
\]
Then the exact matrix in the normalization used by the current
continuations is
\begin{equation}
 (M_{a,b})_{tk}=\mu_{t+k-a+1}
 \left\{
 \frac{(2b+2t-2a+1)!}{(2b-2k)!}
 +\mathbf1_{k<a}\frac{(2t+1)!}{(2a-2k)!}
 \right\},
 \quad 0\le t,k<b.
 \label{eq:source-M}
\end{equation}
Every factorial argument in \eqref{eq:source-M} is nonnegative in the
range where its term is present.  Negative Bernoulli indices are zero
by convention.

The first $b$ coordinates $C_0,\ldots,C_{b-1}$ are free formal labels,
while \eqref{eq:late-coordinates} describes the next $b$.  Hence the
coordinate matrix of $C_0,\ldots,C_{2b-1}$ is block diagonal with
blocks $I_b$ and $M_{a,b}$.  Since the transform
\eqref{eq:C-transform} is triangular, we obtain the fundamental
criterion.

\begin{proposition}[Source rank criterion; inherited]
\label{prop:rank-criterion}
The first $q=2b$ centered moments $D_0,\ldots,D_{2b-1}$ have full
formal rank $2b$ if and only if
\[
 \det M_{a,b}\ne0.
\]
The statement is entirely rational and makes no claim of numerical
linear independence among the values assigned to the coordinate
labels.
\end{proposition}

\subsection{The shifted-Hankel reduction}

The Bernoulli core of \eqref{eq:source-M} is the shifted Hankel matrix
\begin{equation}
 H_{a,b}=(\mu_{t+k-a+1})_{0\le t,k<b}.
 \label{eq:H}
\end{equation}
The source factorization has the form
\[
 M_{a,b}=R(H_{a,b}+U)E,
\]
where $R$ and $E$ are nonsingular diagonal matrices and $U$ has rank at
most $a$.  Orthogonal-polynomial row operations based on the Pad\'e
approximants to
\[
 \frac{\sqrt z}{2}\coth\frac{\sqrt z}{2}
   =\sum_{j\ge0}\mu_jz^j
\]
annihilate the last $b-a$ columns in the final $a$ transformed rows.
Consequently, the full determinant reduces to a retained determinant of
size $a$.  The reduction also proves that the unperturbed matrix
$H_{a,b}$ is nonsingular and gives its inertia~\cite{BTD,QSR,DilcherJiu2022,DilcherJiu2023}.

For use in the order-ten proof, we state the retained matrix explicitly.
Assume $b\ge2a$.  For $0\le i,\ell<a$, set
\begin{align}
 n_i&=b-a+i,&
 N_{i\ell}&=2(n_i-\ell)+1,&
 R_\ell&=2b-2a+2+2\ell,\label{eq:retained-indices}\\
 F_{i\ell}&=\frac1{(2\ell+1)!}
 \prod_{r=0}^{2\ell-1}(2n_i-r)(2n_i+3+r),
 \label{eq:F-entry}
\end{align}
with the empty product equal to one.  Define
\begin{equation}
 W_{i\ell}(z)=(z+2\ell+1)
              \fall{(2n_i+2\ell+z+1)}{2i},
 \qquad
 d_{i\ell h}=\frac1{h!}\sum_{j=0}^h
 (-1)^{h-j}\binom hjW_{i\ell}(N_{i\ell}-j).
 \label{eq:d-finite-difference}
\end{equation}
The $d_{i\ell h}$ are integers because they are Newton coefficients of
an integer-valued polynomial.  With independent symbols
$U_1,\ldots,U_{2a-1}$, put
\begin{equation}
 \A^{[a]}_{i\ell}
 =2F_{i\ell}
 +\frac{R_\ell!}
 {(2n_i+1)(2n_i+2)(2\ell+2)!N_{i\ell}!}
 \sum_{\substack{1\le h\le2i+1\\h\ {\rm odd}}}
 d_{i\ell h}\fall{N_{i\ell}}h
 U_{a-i+\ell+(h-1)/2}.
 \label{eq:retained-A}
\end{equation}
At the Bernoulli specialization
\begin{equation}
 U_j=2B_{2b-2j},\qquad 1\le j\le2a-1,
 \label{eq:U-specialization}
\end{equation}
the retained determinant controls the original source matrix.

\begin{lemma}[Exact determinant multiplier; inherited]
\label{lem:multiplier}
Let
\[
 h_n=\frac{((2n+1)!)^2}{(4n+2)!(4n+3)!},
 \qquad
 c_n=\frac{(2n+1)(2n+2)!}{(4n+2)!},
 \qquad
 \epsilon_{a,b}=a(b-a)+\frac{a(a-1)}2.
\]
For $b\ge2a$,
\begin{align}
 \det M_{a,b}
 ={}&(-1)^{\epsilon_{a,b}}
 \frac{\prod_{t=0}^{b-1}(2b+2t-2a+1)!}
      {\prod_{k=0}^{b-1}(2b-2k)!}
 \left(\prod_{n=0}^{b-a-1}h_n\right)
 \left(\prod_{i=0}^{a-1}c_{b-a+i}\right)
 \det\A^{[a]}.
 \label{eq:exact-multiplier}
\end{align}
Every factor outside $\det\A^{[a]}$ is nonzero.  Therefore
\[
 \det M_{a,b}\ne0
 \quad\Longleftrightarrow\quad
 \det\A^{[a]}\ne0.
\]
\end{lemma}

The delivered verifier checks \eqref{eq:exact-multiplier} modulo the
certificate prime at $b=10,11,12,13,14$, comparing the direct source
determinant with the retained determinant times the full multiplier.
These spot checks are not needed once \cref{lem:multiplier} is accepted,
but they guard against normalization errors in the implementation.

\subsection{Eventual nonsingularity}

The retained determinant is a polynomial of total Bernoulli degree at
most $a$ after clearing explicit denominators.  Its top homogeneous
part is a positive Hankel moment determinant after Euler's formula for
$B_{2m}$ is substituted.  The lower Bernoulli degrees are bounded by an
explicit coefficient norm.  Comparing the two produces a terminating
threshold algorithm for each fixed $a$.

\begin{theorem}[Effective eventual rank; inherited]
\label{thm:eventual}
For every fixed $a\ge1$, there is a constructible $B_a$ such that
$\det M_{a,b}\ne0$ for all $b\ge B_a$.  The exact source computation
gives
\[
 B_1=16,
 \quad B_2=45,
 \quad B_3=142,
 \quad B_4=467,
 \quad B_5=1920.
\]
In particular,
\begin{equation}
 \det M_{5,b}\ne0\qquad(b\ge1920).
 \label{eq:eventual-order-ten}
\end{equation}
\end{theorem}

The role of the present paper is therefore finite but decisive: prove
nonsingularity for every $6\le b<1920$ without leaving exceptional
progressions.
